% Generated from validated Article Draft Result. Do not hand-edit; revise the structured draft instead.
\section{AutonomyにはGuardrailが要る — Prompt InjectionとData Exfiltration}
\noindent\textbf{Agentがweb・tool・fileへ接続されると、prompt injectionは単なる誤応答ではなく外部actionとdata exfiltrationの問題になる。1月は、その具体的な防御設計と攻撃面の体系化が同時に現れた。}\autocite{src-0d2abbb0fa62,src-e9010a7fe30b}

OpenAIが1月28日に説明した脅威は、prompt-injected contentがAgentにURLを取得させ、そのURLへconversation由来のsecretを埋め込むことで外部へ情報を運ぶというものだ。Agentが外部contentを読むだけでなく、fetchというactionを実行できることで、prompt injectionがdata flowの問題へ変わる。\autocite{src-0d2abbb0fa62}


説明されたsafe-URL safeguardでは、独立にpublicと観測されたURLは自動取得を許し、それ以外はfetchをblockするかexplicit user actionを要求する。重要なのは、Agentが『読める』ことと『外へ通信してよい』ことの間に確認境界を設ける発想である。\autocite{src-0d2abbb0fa62}


\begin{claimboundary}
この機構はURL-based exfiltrationに対する一層の防御であり、arbitrary web contentの安全性やprompt injection一般の解決を意味しない。さらにsecurity propertyはOpenAIのimplementation descriptionで、本稿がpenetration testで再検証したものではない。\autocite{src-0d2abbb0fa62}
\end{claimboundary}

独立研究側では、Prompt Injection Attacks on Agentic Coding Assistantsがcoding assistantsだけでなくskills、tools、protocol ecosystemsまでを対象にしたSystematization of Knowledgeを提示した。paperのabstractは78 studies、42 attack techniques、18 defense mechanismsの整理を報告している。\autocite{src-e9010a7fe30b}


\begin{claimboundary}
78/42/18という数は異なる研究を束ねたauthor-reported synthesisであり、一つのcontrolled experimentの成功率として読んではならない。また保持しているsourceはJanuaryのarXiv v1であり、後続研究によってlandscapeは更新され得る。\autocite{src-e9010a7fe30b}
\end{claimboundary}

1月のAgent securityが示すのは、modelへのinstruction attackだけを見るのでは足りないということだ。tool、file、skill、protocol、network fetchへ権限が伸びるほど、どのactionを自動実行し、どこでuser confirmationを要求するかがarchitecture上の核心になる。\autocite{src-0d2abbb0fa62,src-e9010a7fe30b}
